\chapter{Discussion}
\label{ch:discussion}

This chapter interprets the findings in Chapter~\ref{ch:results} and develops a framework for selecting an
extraction method and, where needed, a serialization format. The populations follow
Section~\ref{sec:analysis-cohorts}: retrieval accuracy uses the answerable cohort, while comparative
operational efficiency uses the context-fit cohort. These overlapping but different populations support
separate interpretations, not a jointly measured quality--resource frontier.

\section{Interpretation of the main findings}
\label{sec:main-findings-interpretation}

\subsection{Pipeline selection and database size}
\label{sec:pipeline-size-interpretation}

The paired comparisons in Section~\ref{sec:common-cohort-comparison} favour Text-to-SQL for reference-answer
agreement across metrics on the evaluated relational workloads. For example, on Spider, even JSON, the direct
format with the highest mean NEM, had a paired JSON-minus-SQL difference of \num{-0.465}
(\numrange{-0.507}{-0.422}). The corresponding comparisons on sampled BIRD also favoured SQL.
Delegating retrieval and aggregation to the database engine was therefore more effective than asking the
model to derive the answer directly from the serialized records under these conditions. This supports
considering the extraction method before optimizing its input representation.

The size-tier results strengthen this interpretation. In
Section~\ref{sec:pipeline-accuracy-by-database-size}, every paired direct-minus-SQL interval was below zero
for all three metrics in every tier, including small databases. Small size establishes potential
direct-context feasibility, but does not by itself justify preferring Direct Context Inference. Both
pipelines had lower mean NEM in the large tier, and the size of the gap changed differently across metrics.
The evidence does not support a simple rule that one method becomes preferable at a particular tier
boundary.

The negative database-size associations in Section~\ref{sec:direct-results} include zero scores for context
failures, rather than isolating accuracy among feasible requests. Size also varied with domain, schema, and
questions. In the pooled tiers, all small and medium databases
were from Spider, while all sampled BIRD databases belonged to the large tier. Dataset composition is
therefore an additional confounder when comparing tiers. Neither these comparisons nor the separate
within-dataset associations establish a causal effect of size. They instead show why context capacity and
answer quality must be considered separately. For relational tasks that can be expressed in SQL and whose
infrastructure permits its use, the results support Text-to-SQL as the first candidate; the deployment
conditions are specified in Section~\ref{sec:decision-framework}.

\subsection{Serialization trade-offs}

CSV is a practical starting representation for flat Direct Context Inference inputs. It required the
fewest payload tokens on average in Section~\ref{sec:serialization-results}, while TOON remained close.
The request-level comparison in Section~\ref{sec:direct-results} also showed lower input-token use and
shorter mean latency for these compact formats than for the more verbose alternatives. Much of the timing
difference occurred in TTFT, so the benefit primarily concerned processing the supplied input rather than
shortening the generated answer. Compactness also extended the set of databases that could potentially
be supplied to the model.

However, token efficiency does not establish a universal agreement ranking. JSON had the highest observed
direct NEM mean on Spider, but its paired advantage over CSV was only \num{0.57} percentage points
(\numrange{-1.21}{2.34}). This does not establish a clear advantage, nor does it prove equivalence. CSV
had the highest direct mean BERTScore and ROUGE-L on Spider and the highest means on all three measures
on sampled BIRD, with TOON generally competitive. Format selection should therefore consider the paired
evidence and its uncertainty, not only the ordering of means.

These findings favour avoiding unnecessary verbosity, rather than conversion at any price. Representational
fidelity remains a prerequisite: reducing tokens is not useful if it removes relationships needed for the
answer. The experiment evaluated serialized relational tables, not nested-data retrieval. Moreover,
payload counts describe the data alone, not a complete request, and arithmetic means can be dominated by
one large database. Actual context feasibility must be checked for the selected representation and model.
The framework gives the conditional rules for preserving structure or retaining an existing format.

\subsection{Model and dataset effects}

Model selection mattered within both pipelines, but nominal parameter count was not an isolated explanatory
variable. The family-specific comparisons in Figures~\ref{fig:direct-accuracy-by-model-size} and
\ref{fig:text-to-sql-accuracy-by-model-size} showed different gains across variants and datasets. In
Text-to-SQL, the Qwen comparisons showed substantial improvements from the smallest variants, with much
of the observed gain occurring by 4B. This does not imply that the same scaling pattern holds for every
family or application. In Direct Context Inference, the larger Gemma and Ministral variants did not show
clear NEM improvements in the reported sampled-BIRD comparisons.

Architecture and training also changed between variants. Gemma~4 26B-A4B, for example, is positioned at
its total nominal parameter count in the figures, not its approximately 4B active parameters. Model choices
therefore require measured performance rather than a parameter-count threshold.

Lower agreement on sampled BIRD in both pipelines further shows that the workload matters. Its data,
schemas, questions, and sampling procedure differ from Spider, so the difference cannot be attributed to
one unisolated cause, such as data cleanliness or schema size. The size comparisons discussed above provide
another view of this dependence, not a separate explanation. A configuration that works on one dataset
must still be evaluated on the intended questions and data.

\subsection{Operational implications}

The shared context-fit comparison provides evidence about resource use without restricting inclusion to
successful runs. It retained SQL failures and all recorded correction work. Text-to-SQL used fewer mean
input tokens than every direct format (Figure~\ref{fig:common-cohort-token-usage}), consistent with
supplying a schema instead of complete database contents. Output-token differences were less consistent:
the Spider CSV-to-SQL comparison did not establish a clear difference. Smaller input does not necessarily
mean a shorter completion, especially when reasoning and correction outputs are included.

Mean total latency also favoured Text-to-SQL (Figure~\ref{fig:common-cohort-latency}). Even for CSV, the
CSV-to-SQL ratio was \num{2.69} (\numrange{1.06}{4.27}) on Spider. On the two sampled BIRD databases in
the context-fit cohort, it was \num{10.43} (\numrange{9.54}{11.51}). The BIRD estimate must remain
limited to those two databases. Most of the latency difference occurred in TTFT, while SQL execution
contributed little under the evaluated conditions. This explains why adding an execution step did not
cancel the benefit of the smaller prompt. An aggregate advantage nevertheless does not guarantee faster
responses for every run or model.

The separate answerable-cohort Text-to-SQL results also caution against equating smaller models with less
total work. Qwen~3.5 2B used fewer input and output tokens and had lower mean complete-run latency than
0.8B, alongside fewer retries (Section~\ref{sec:text-to-sql-results}). This concerns the whole attempt
sequence, not faster individual generations by the larger model. Retry-specific implications are discussed
in Section~\ref{sec:sql-error-correction}.

For paid APIs, lower input-token consumption could reduce expenditure under the relevant pricing policy;
it is not a measured monetary saving here. Production SQL execution could also be more expensive than in
these experiments. Likewise, concurrency, throughput, and prompt caching could change the operational
balance, but their effects were not measured. The local latency results support a starting choice, not a
production performance guarantee.

\section{Relationship to prior work}

The serialization findings are consistent with research treating representation as both an efficiency and
an interpretation concern. StrucText-Eval \cite{struct-text-eval} shows that agreement depends on the
format and model, while ONTO and TOON emphasize reducing token overhead
\cite{onto-serialization-format,toon-format-2025}. Here, CSV remained the most compact representation on
average and TOON was competitive, rather than uniformly superior. This is not a contradiction: the
evaluated inputs were relational tables, and efficiency was measured against several existing formats.
Retrieval scores also included context failures. Claims about one format therefore depend on the task,
serializer, tokenizer, and comparison population, not only on the format's intended compactness.

Long-context research shows that fitting information into a context window does not ensure that a model
uses it reliably \cite{liu2024lost}. The present results are compatible with that distinction: database
size was associated with lower direct-context scores, and small inputs did not remove the observed SQL
advantage. However, we did not vary evidence position or measure lost-in-the-middle behaviour. Context
failures and differences between databases also contributed to the observed associations. The practical
implication is to evaluate retrieval quality after checking feasibility, rather than treating the advertised
context capacity as a guarantee of useful extraction.

The minimalist architecture in Section~\ref{sec:minimalist-text-to-sql} establishes a baseline, not the
limit of Text-to-SQL performance. Metadata enrichment \cite{shkapenyuk2025metadata}, multi-candidate
generation and selection \cite{wang2025agentar}, and targeted verification \cite{li2026deepeyesql}
address information or reasoning gaps beyond execution-error repair. These extensions also add profiling,
retrieval, generation, and selection work whose preprocessing and repeated-use costs require separate
measurement. Their published execution-accuracy results cannot be ranked directly against this thesis's
textual agreement or execution-success results, particularly on sampled BIRD. They instead motivate
evaluating whether additional workflow complexity provides sufficient application benefit to justify its
resource requirements.

\section{SQL error correction: benefits and limitations}
\label{sec:sql-error-correction}

Under the policy in Section~\ref{sec:minimalist-text-to-sql}, correction increased execution success on
the answerable cohort. Table~\ref{tab:text-to-sql-error-correction} shows that most recoveries occurred
on the first retry, with fewer additional recoveries from the remaining attempts. Database feedback was
therefore useful for some initially failed queries, but repeatedly supplying feedback did not reliably
resolve the remaining failures. As distinguished in Section~\ref{sec:sql-execution-metrics}, recovery
means obtaining executable SQL, not necessarily recovering the intended answer. An executable query with
the wrong join or filter does not trigger this correction policy.

Terminal failures consumed \qty{95.3}{\percent} of recorded retry-generation time on Spider and
\qty{88.3}{\percent} on sampled BIRD. These runs exhausted every available retry, whereas recovered runs
stopped as soon as a query executed. This concentration reflects both persistent failure and the stopping
rule. Here, retry overhead is model-generation time, including retry
TTFT. Per-attempt SQL execution time was unavailable, so it is not a complete end-to-end measure of
correction latency.

A retry budget should consequently reflect the value of incremental execution recovery and the
application's latency and resource requirements. One retry is a reasonable starting policy for a
latency-sensitive application because it captured most observed recoveries. It is not an established
optimum: the study did not optimize an application objective or evaluate adaptive stopping policies.
Additional retries may be justified where the remaining recoveries matter, provided that answer quality
is validated separately and the added work fits the application budget.

\section[Decision framework]{Decision framework for structured data information extraction}
\label{sec:decision-framework}

\subsection{Decision inputs}

The framework is evidence-informed guidance, not an independently validated optimal selection procedure.
Its inputs are the source representation and relationships that must be preserved;
SQL applicability, engine availability, and access permissions; the context budgets and capabilities of
available models; and the application's answer-quality, latency, and resource requirements. These inputs
determine the candidate configurations. No universal acceptance thresholds follow from the benchmark results.

\subsection{Selection rules}

Text-to-SQL is the first candidate when the data have a faithful relational representation, the tasks are
SQL-expressible, an appropriate engine and access are available, and the complete schema request fits
with sufficient output allowance. Section~\ref{sec:pipeline-size-interpretation} provides the empirical
rationale; small database size alone does not reverse that starting preference. The model and retry budget
remain configuration choices, with the latter informed by Section~\ref{sec:sql-error-correction}.
Inspectable SQL can assist auditing, but does not establish inherently correct or reproducible outputs.

Direct Context Inference remains a candidate when SQL cannot be used, or when a tested SQL configuration
does not meet requirements and a feasible direct alternative exists. CSV is the starting format for flat
inputs. Nested inputs require a structure-preserving representation; TOON is a candidate, not a demonstrated
retrieval winner for nested data. An existing format can be retained when an interface requires it or when
conversion offers little practical benefit. Conversion time was not measured in this study.

For either pipeline, feasibility requires checking the complete request, including instructions, the
question, and chat-template overhead, together with an allowance for output and any reasoning tokens.
Direct representation selection must precede this check because format and tokenizer affect the result.
This is a deployment recommendation, distinct from the recorded experimental context-fit criterion in
Section~\ref{sec:analysis-cohorts}.

\subsection{Decision tree and application validation}

\begin{figure}[p]
  \centering
  \resizebox{\textwidth}{!}{\begingroup
\renewcommand{\baselinestretch}{1}
\fontsize{10}{11.5}\selectfont
% Natural size: about 13.6 cm by 16.6 cm; use \resizebox{\textwidth}{!}{...}.
\begin{tikzpicture}[
    x=1cm, y=1cm,
    font=\fontsize{10}{11.5}\selectfont,
    every node/.style={align=center},
    box/.style={draw, rounded corners=2pt, inner sep=5pt, text width=4.8cm},
    gate/.style={box, fill=black!6},
    outcome/.style={box, line width=0.8pt},
    arrow/.style={-{Stealth[length=2mm]}, line width=0.6pt},
    answer/.style={fill=white, inner sep=1pt}
  ]
  \node[gate, text width=10.2cm, minimum height=2.4cm] (sqlgate) at (0,0) {
    \textbf{SQL prerequisites met?}\\[2pt]
    Faithful relational representation\\
    SQL-expressible workload\\
    SQL engine and access available\\
    Complete schema request + output allowance fits
  };
  \node[box, minimum height=0.95cm] (sql) at (-3.1,-3.2) {
    \textbf{Text-to-SQL candidate}
  };
  \node[box, minimum height=3.2cm] (format) at (3.1,-3.5) {
    \textbf{Direct Context Inference}\\
    \textbf{Select representation}\\[2pt]
    Flat: CSV starting choice\\
    Nested: preserve structure;\\
    TOON is a candidate\\[2pt]
    Retain existing representation\\
    when appropriate
  };
  \node[gate, minimum height=1.2cm] (directfit) at (3.1,-6.3) {
    \textbf{Complete direct request}\\
    \textbf{+ output allowance fits?}
  };
  \node[gate, text width=9.6cm, minimum height=2.15cm] (accept) at (0,-8.65) {
    \textbf{Application requirements met?}\\[2pt]
    Answer quality on representative tasks\\
    SQL execution reliability (if applicable)\\
    Latency and resource use
  };
  \node[outcome, minimum height=0.9cm] (select) at (-3.1,-10.7) {
    \textbf{Select configuration}\\
    Method, model, representation
  };
  \node[gate, minimum height=2.45cm] (reassess) at (3.3,-12.7) {
    \textbf{Reassess candidates}\\[2pt]
    New model or configuration\\
    or another method\\
    Direct option remains open\\[2pt]
    \textbf{Untested viable candidate?}
  };
  \node[gate, minimum height=2.45cm] (recheck) at (-3.3,-12.7) {
    \textbf{Repeat relevant checks}\\[2pt]
    SQL prerequisites, or direct\\
    representation selection;\\
    complete request and\\
    output allowance fit
  };
  \node[outcome, minimum height=0.9cm] (other) at (3.3,-14.9) {
    \textbf{Another approach needed}
  };

  \draw[arrow] (sqlgate.south -| sql.north) -- node[answer, midway] {Yes} (sql.north);
  \draw[arrow] (sqlgate.south -| format.north) -- node[answer, midway] {No} (format.north);
  \draw[arrow] (format.south) -- (directfit.north);
  \draw[arrow] (sql.south) -- (accept.north -| sql.south);
  \draw[arrow] (directfit.south) -- node[answer, midway] {Yes} (accept.north -| directfit.south);
  \draw[arrow] (directfit.east) -- node[answer, midway] {No}
    (6.8,-6.3) -- (6.8,-12.7) -- (reassess.east);
  \draw[arrow] (accept.south -| select.north) -- node[answer, midway] {Pass} (select.north);
  \draw[arrow] (accept.south -| reassess.north) -- node[answer, midway] {Fail} (reassess.north);
  \draw[arrow] ([yshift=4mm]reassess.west) -- node[answer, midway] {Yes}
    ([yshift=4mm]recheck.east);
  \draw[arrow] ([yshift=-4mm]recheck.east) -- node[answer, midway] {Unfit}
    ([yshift=-4mm]reassess.west);
  \draw[arrow] (recheck.west) -- (-6.8,-12.7) -- (-6.8,-8.65)
    -- node[answer, midway] {Fits} (accept.west);
  \draw[arrow] (reassess.south) -- node[answer, midway] {No} (other.north);
\end{tikzpicture}
\endgroup
}
  \caption{Evidence-informed selection of an extraction method and representation. Feasible candidates
    share application checks for answer quality, SQL execution reliability where applicable, latency,
    and resource use. Failure leads to reassessment, including a direct alternative after SQL failure.
    TOON is a structure-preserving candidate; nested-data retrieval was not tested.}
  \label{fig:decision-framework}
\end{figure}

Figure~\ref{fig:decision-framework} presents the decision sequence. Both pipelines must be evaluated on
representative application questions, separating answer quality from SQL execution reliability where
applicable and measuring latency and resource use on those same questions. Passing execution checks alone
is insufficient. Where several candidates meet the requirements, application priorities determine the
selection.

Failure of a feasibility or validation check calls for reassessing available models, configurations, or
methods, rather than automatically selecting the other pipeline. SQL validation failure does not rule out
a feasible direct alternative. Any revised candidate must pass the relevant checks again. If none can
meet the requirements, another approach is needed, potentially involving selective retrieval or a different
extraction workflow. Such extensions were not evaluated here.

\section{Threats to validity and study limitations}
\label{sec:validity-limitations}

\subsection{Measurement and task representation}

The metric limitations in Section~\ref{sec:reference-answer-metrics} affect comparisons between free-text
direct answers and rendered SQL rows. Ordering, explanation, duplicates, and numeric formatting can change
agreement without a corresponding change in factual correctness. The three measures do not remove this
shared dependence on text. The benchmarks also consist of
SQL-expressible tasks. The results therefore characterize this task setting, not every form of information
extraction from structured data.

\subsection{Sampling and population coverage}

Independent BIRD table sampling, described in Chapter~\ref{ch:implementation}, can remove related rows,
change joins and aggregates, and alter task difficulty even when the reference remains non-empty.
Recomputing references ensures that both pipelines use the same experimental data; excluding empty
references does not restore the original tasks. Conclusions about sampled BIRD are consequently not
equivalent to results on its full databases. The context-fit restriction further narrows operational
coverage, especially on BIRD, where only two databases remained. Narrow bootstrap intervals over those
databases do not imply broad precision across BIRD workloads. The distinction between feasibility and
size-related agreement is discussed in Section~\ref{sec:pipeline-size-interpretation}.

\subsection{Experimental controls and uncertainty}

The results apply to the serializers, prompts, retained model variants, quantization, context budget, and
local serving configuration in Chapters~\ref{ch:methodolgy} and \ref{ch:implementation}. Different
formatting or prompting choices could change the comparisons, and parameter count was not independently
isolated from architecture and training. The exclusion of Qwen~3.5 27B also limits the evaluated scaling
range. The intervals in Section~\ref{sec:confidence-intervals} describe pointwise uncertainty, not a
simultaneous guarantee across the many format and model comparisons. Neither the highest observed mean
nor an interval containing zero establishes a universal ranking or equivalence.

\subsection{Deployment and framework validity}

The framework has not been independently evaluated on new workloads. No monetary expenditure,
production-scale SQL execution, concurrency effects, or caching benefit was measured. Recorded latency
excludes the preprocessing specified in Section~\ref{sec:evaluation-metrics-implementation}; production
decisions may need to include that work and its amortization across requests. The local retry-time
qualification is given in Section~\ref{sec:sql-error-correction}. Nested and non-relational retrieval
were not tested, so the framework's corresponding branches identify candidates for validation rather than
established performance advantages. The final chapter states the thesis's contribution within these
boundaries and identifies the next investigations.
